%% kapa: exam=1 round=28 year=2017 period=1 form=A subject=civil_law subject_ko=민법 seq=2 number=2
%% source: Q-Net 감정평가사 기출문제 (cst00309) artlSeq=5199603 file=extracted/1008/5199603/2017년 제28회 감정평가사 1차시험 문제지/1교시 A형.pdf
%% conversion: kordoc 4.24.51 -> tools/gen1.py (auto); review=proofread-sonnet; flags=none
\begin{question}{2}
능력에 관한 설명으로 옳은 것은? (다툼이 있으면 판례에 따름)
\begin{choices}
  \item 2인 이상이 동일한 위난으로 사망한 경우에는 동시에 사망한 것으로 본다.
  \item 태아는 불법행위로 인한 손해배상청구권에 관하여 이미 출생한 것으로 추정한다.
  \item 태아는 그 법정대리인에 의하여 수증행위를 할 수 있다.
  \item 제한능력을 이유로 법률행위를 취소한 경우, 제한능력자는 선의ㆍ악의를 묻지 아니하고 그 행위로 인하여 받은 이익이 현존하는 한도에서 상환할 책임이 있다.
  \item 계약자유의 원칙상 제한능력자를 보호하는 규정에 반하는 매매계약도 유효하다.
\end{choices}
\end{question}
\answer{4}
